\begin{ouverture}[calcul]{10}{Calculer avec les décimaux}{Partie 2 · Fractions et nombres décimaux}
{Sommes · Différences · Produits par un entier}
\objectif{Poser des additions et soustractions de décimaux.}
\objectif{Multiplier un décimal par 10 et par un entier inférieur à 10.}
\objectif{Résoudre des problèmes de mesures et de monnaie.}
\prerequis{Rangs décimaux et opérations posées sur les entiers.}
\end{ouverture}
\pagestyle{otMain}\essentiel\fontsize{11}{15}\selectfont
\phantomsection\label{cours-10-1}
\essSection{1}{Additionner et soustraire}
\begin{savoirfaire}[Aligner les virgules]Poser une addition ou une soustraction de décimaux, c’est aligner les virgules : unités sous unités, dixièmes sous dixièmes. Compléter les rangs vides par des zéros : pour $10-3{,}48$, écrire $10{,}00$.\end{savoirfaire}
\begin{exemple}$4{,}7+2{,}85$ : on pose $4{,}70+2{,}85$. $0+5$ centièmes, $7+8$ dixièmes font 15 dixièmes : on écrit 5 et on retient une unité. Résultat : $7{,}55$. Estimation : proche de $5+3=8$.\end{exemple}
\begin{propriete}Les techniques des entiers s’appliquent rang par rang : dix centièmes s’échangent contre un dixième, dix dixièmes contre une unité. La virgule du résultat s’aligne sous les autres.\end{propriete}
\begin{savoirfaire}[Calculer de tête]De tête : ajouter un entier ne touche que la partie entière ($5{,}34+2=7{,}34$). Chercher les compléments : $0{,}35+0{,}65=1$ et $7{,}2+2{,}8=10$. Pour rendre la monnaie sur 10 euros après $7{,}45$ : $7{,}45\to7{,}50\to8\to10$, soit $2{,}55$.\end{savoirfaire}
\clearpage
\phantomsection\label{cours-10-2}
\essSection{2}{Multiplier un décimal}
\begin{propriete}Multiplier par 10 fait monter chaque chiffre d’un rang : $3{,}7\times10=37$ et $0{,}85\times10=8{,}5$. Diviser par 10 fait l’inverse : $45{,}2:10=4{,}52$. La virgule semble glisser : ce sont les chiffres qui changent de rang.\end{propriete}
\begin{exemple}$2{,}4\times3$ : 24 dixièmes multipliés par 3 font 72 dixièmes, soit $7{,}2$. Pour $3{,}45\times6$ : $345$ centièmes $\times\,6=2\,070$ centièmes, soit $20{,}70=20{,}7$. Estimation : $3{,}5\times6=21$.\end{exemple}
\begin{savoirfaire}[Poser un produit par un entier]Poser $3{,}45\times6$ comme $345\times6$, puis replacer la virgule : le résultat garde deux chiffres après la virgule, comme $3{,}45$. Toujours contrôler avec un ordre de grandeur.\end{savoirfaire}
\begin{attention}Multiplier par 10 n’ajoute pas un zéro à un décimal : $3{,}7\times10$ ne fait pas « $3{,}70$ » mais 37. Penser aux rangs : chaque chiffre devient dix fois plus grand.\end{attention}

\clearpage\rubrique{Découverte}
\begin{blocLivre}{ex-10-1}\begin{exercice}[1]{Une somme posée}
Poser et calculer $4{,}7+2{,}85$.
\end{exercice}\end{blocLivre}
\begin{blocLivre}{ex-10-2}\begin{exercice}[1]{Une différence posée}
Poser et calculer $10-3{,}48$.
\end{exercice}\end{blocLivre}
\begin{blocLivre}{ex-10-3}\begin{exercice}[1]{De tête, sans retenue}
Calculer mentalement $5{,}34+2$, puis $5{,}34+0{,}2$, puis $7{,}89-4$.
\end{exercice}\end{blocLivre}
\begin{blocLivre}{ex-10-4}\begin{exercice}[1]{Par 10, dans les deux sens}
Calculer $3{,}7\times10$ ; $0{,}85\times10$ ; $45{,}2:10$.
\end{exercice}\end{blocLivre}
\begin{blocLivre}{ex-10-5}\begin{exercice}[1]{Un produit en dixièmes}
Calculer $2{,}4\times3$ en raisonnant en dixièmes.
\end{exercice}\end{blocLivre}

\rubrique{Entraînement}
\begin{blocLivre}{ex-10-6}\begin{exercice}[2]{Compléments utiles}
Compléter : $0{,}35+\square=1$, puis $7{,}2+\square=10$.
\end{exercice}\end{blocLivre}
\begin{blocLivre}{ex-10-7}\begin{exercice}[2]{Un produit posé}
Poser et calculer $3{,}45\times6$.
\end{exercice}\end{blocLivre}
\begin{blocLivre}{ex-10-8}\begin{exercice}[2]{Regrouper pour simplifier}
Calculer $2{,}6+4{,}85+0{,}4$ en regroupant deux termes bien choisis.
\end{exercice}\end{blocLivre}
\begin{blocLivre}{ex-10-9}\begin{exercice}[2]{Contrôler une soustraction}
Noé calcule $8-2{,}35=6{,}35$. Vérifier avec une addition et corriger si besoin.
\end{exercice}\end{blocLivre}
\begin{blocLivre}{ex-10-10}\begin{exercice}[2]{Rendre la monnaie}
Un achat coûte $7{,}45$ euros, payé avec un billet de 10 euros. Calculer la monnaie rendue en passant par des étapes simples.
\end{exercice}\end{blocLivre}

\rubrique{Problèmes}
\begin{blocLivre}{ex-10-11}\begin{exercice}[3]{Le panier de courses}
Une baguette coûte $1{,}20$ euro, une plaquette de beurre $2{,}85$ euros et une brique de lait $1{,}15$ euro.\begin{questions}
\item Calculer la dépense totale.
\item Quelle monnaie rend-on sur un billet de 10 euros ?
\end{questions}
\end{exercice}\end{blocLivre}
\begin{blocLivre}{ex-10-12}\begin{exercice}[3]{Les sorties du coureur}
Un coureur parcourt $4{,}5$ km lundi, puis $1{,}8$ km de moins mardi. Quelle distance parcourt-il mardi ? Et sur les deux jours ?
\end{exercice}\end{blocLivre}
\begin{blocLivre}{ex-10-13}\begin{exercice}[3]{Trois carnets identiques}
Un carnet coûte $2{,}45$ euros. Calculer le prix de trois carnets, puis la monnaie rendue sur 10 euros.
\end{exercice}\end{blocLivre}
\begin{blocLivre}{ex-10-14}\begin{exercice}[3]{Le ruban coupé}
Dans un ruban de $2{,}75$ m, on coupe deux morceaux de $0{,}9$ m chacun. Quelle longueur de ruban reste-t-il ?
\end{exercice}\end{blocLivre}
\begin{blocLivre}{ex-10-15}\begin{exercice}[3]{L’alignement fautif}
Tom pose $5{,}6+2{,}14$ en alignant les nombres à droite, comme des entiers, et trouve $2{,}70$. Expliquer son erreur et donner le bon résultat.
\end{exercice}\end{blocLivre}

\clearpage
